hl|tc| Planning and Learning with Tabular Methods
In addition to the direct reinforcement learning approaches which improve value/policy with experience gained by acting on the environment, a model of the environment can be used to improve the process.
Model based and Direct learning can also be viewed in a unified manner : for each iteration of direct RL, there can be \(n \) iterations of model learning,
li| Value / Policy \( \rightarrow \) Actions \( \rightarrow \) Experience \( \rightarrow \) Direct Learning \( \rightarrow \) Value / Policy
li| Value / Policy \( \rightarrow \) Actions \( \rightarrow \) Experience \( \rightarrow \) Model Learning \( \rightarrow \) Model \( \rightarrow \) Planning \( \rightarrow \) Value / Policy
Dyna-Q is one such algorithm, which is basically a Tabular Dyna a combination of Direct Q-Learning, and intermittent planning steps based on a model.
When models are incorrect due to limited samples or simply because the enviroment has changed, the suboptimal policy is quickly corrected if the model is optimistic. If the environment becomes better through some paths, optimal policy may never be found without some incentive added to the reward for exploration i.e. \( r \leftarrow r + \kappa \sqrt{\tau} \) where \( \kappa \) is a small constant and \( \tau \) is the count of steps since a transition was explored.
Prioritized Sweeping is a common optimization technique that can dramatically speed up the basic RL models. The general idea with this is backward focusing of updates. A priority queue keyed by the difference \( R + \gamma \max_a Q(S', A) - Q(S, A) \) can be used to update the Q-function for every step. Note that covering new elements to be added to the priority queue requires a model for predicting state-action pairs that lead to the current pair being iterated on.
In stochastic environments with large branching factors, sample updates can be significantly better to save computation time because the number of steps required for sampling error to diminish will be much smaller than the number of steps required to sweep entire state space.
Trajectory Sampling is another technique that can allow the agent to focus only on the relevant states, especially in cases where there are a limited number of starting states. Real-time Dynamic Programming is one such approach that can be much faster than conventional sweep based approaches.
One key property of all of the above algorithms is that they perform planning in the background. The other category of learning algorithms involve decision-time planning. These are particularly effective when working with environments where storing all the states in not feasible, helpful or desirable.
The classical Heuristic Search algorithms to play chess & go used planning at decision time : the agent will traverse the tree of possible trajectories starting at the current state, and apply (possibly hardcoded) heuristic function to estimate leaf values before backing up the estimated values.
Rollout algorithms are decision-time planning algorithms based on Monte Carlo control applied to simulated trajectories that all start at the current state. The goal here is not to find the optimal policy, but to improve upon the rollout policy by averaging returns of sample trajectories.
Monte Carlo Tree Search (MCTS) is a strikingly successful example of decision-time planning. MCTS consists of the following steps,
li| Selection : Starting at the current state, follow a tree policy to reach leaf nodes on the basis of action values attached to the edges of the tree. For example, this policy could select actions based \( \epsilon \)-greedy or UCB selection rule to balance exploration and exploitation.
li| Expansion : On some actions, add one or more child nodes from the selected leaf via unexplored actions.
li| Simulation : Follow the rollout policy to perform simulations starting at the [possibly new] selected leaf.
li| Backup : Updated estimates on the traversed path with the returns observed from the full Monte Carlo trial.
By incrementally expanding and trimming the tree, MCTS effectively grows a lookup table to store partial state-action value function, with focus on high-yielding sample trajectories. MCTS thus retains the benefit of using past experience to guide exploration, but avoids the problem of globally approximating value function.